% ===== PRÉAMBULE CANONIQUE — à coller VERBATIM en tête de CHAQUE .tex =====
\documentclass[11pt,a4paper]{article}
\usepackage[utf8]{inputenc}\usepackage[T1]{fontenc}\usepackage{lmodern}
\usepackage[french]{babel}
\usepackage{amsmath,amssymb,mathtools}\usepackage{icomma}
\usepackage[version=4]{mhchem}          % \ce{...} : équations & formules
\usepackage{siunitx}\sisetup{locale=FR} % \qty{}{}, \si{}, \ang{}
\usepackage{chemfig}                    % \chemfig{...} : structures développées
\usepackage{xcolor}\usepackage{colortbl}
\usepackage{tikz}\usepackage{pgfplots}\pgfplotsset{compat=1.18}
\usepackage[most]{tcolorbox}\usepackage{enumitem}\usepackage{array,booktabs}
\usepackage[a4paper,margin=2.2cm]{geometry}
\usetikzlibrary{arrows.meta,calc,angles,quotes,positioning,patterns,backgrounds,shapes.geometric}
\definecolor{primary}{RGB}{222,49,99}\definecolor{primarydark}{RGB}{150,22,55}
\definecolor{defcol}{RGB}{0,114,178}\definecolor{methcol}{RGB}{52,92,148}
\definecolor{excol}{RGB}{0,135,90}\definecolor{attncol}{RGB}{200,40,55}
\tcbset{boxbase/.style={enhanced,breakable,boxrule=0.9pt,arc=2.5pt,
  left=3mm,right=3mm,top=2.3mm,bottom=2.3mm,fonttitle=\bfseries,coltitle=white}}
% --- boîtes TUTORIEL (noms EXACTS, ne pas renommer) ---
\newtcolorbox{cle}[1][]{boxbase,colback=defcol!6,colframe=defcol,
  title={\textsc{À retenir}\ifx\relax#1\relax\relax\else\textmd{ : \itshape #1}\fi}}
\newtcolorbox{methode}[1][]{boxbase,colback=methcol!7,colframe=methcol,sharp corners,
  title={\textsc{Méthode}\ifx\relax#1\relax\relax\else\textmd{ --- \itshape #1}\fi}}
\newtcolorbox{exemplebox}[1][]{boxbase,colback=excol!5,colframe=excol,
  title={\textsc{Exemple}\ifx\relax#1\relax\relax\else\textmd{ : \itshape #1}\fi}}
\newtcolorbox{piege}[1][]{boxbase,colback=attncol!6,colframe=attncol,
  title={\textsc{Piège}\ifx\relax#1\relax\relax\else\textmd{ : \itshape #1}\fi}}
\newtcolorbox{remarque}[1][]{boxbase,colback=black!4,colframe=black!45,
  title={\textsc{Remarque}\ifx\relax#1\relax\relax\else\textmd{ : \itshape #1}\fi}}
% --- boîtes EXERCICE / CORRIGÉ (noms EXACTS, auto-comptées) ---
\newtcolorbox[auto counter]{exo}[1][]{boxbase,colback=primary!4,colframe=primary,
  title={\textsc{Exercice}~\thetcbcounter\ifx\relax#1\relax\relax\else\textmd{ --- \itshape #1}\fi}}
\newtcolorbox[auto counter]{corrige}[1][]{boxbase,colback=excol!4,colframe=excol,
  title={\textsc{Corrigé}~\thetcbcounter\ifx\relax#1\relax\relax\else\textmd{ --- \itshape #1}\fi}}
\newcommand{\R}{\mathbb{R}}\newcommand{\N}{\mathbb{N}}\newcommand{\Z}{\mathbb{Z}}\newcommand{\ds}{\displaystyle}
% (un \usepackage{fancyhdr}+en-tête est possible ; il est ignoré côté web)
% =========================================================================
\begin{document}

\begin{center}{\large\bfseries\color{primarydark} Thème 5 — Niveau Moyen}\end{center}
\emph{Données (\si{\gram\per\mole}) :} \ce{H2O}$=18$, glucose \ce{C6H12O6}$=180$, \ce{NaOH}$=40$.

\begin{exo}[préparer une solution diluée]
On dispose d'une solution mère de \ce{NaOH} à $C_1 = \SI{5.0}{\mole\per\litre}$. On veut
$\SI{100}{\milli\litre}$ de solution à $C_2 = \SI{0.25}{\mole\per\litre}$.
\begin{enumerate}[label=\alph*)]
  \item Quel volume $V_1$ de mère faut-il prélever~?
  \item Décris brièvement le protocole.
\end{enumerate}
\end{exo}

\begin{exo}[dilution impossible]
On dispose d'une solution à $C_1 = \SI{2.0}{\mole\per\litre}$. Peut-on, \emph{par dilution}, préparer
$\SI{100}{\milli\litre}$ d'une solution à $\SI{4.0}{\mole\per\litre}$~? Justifie.
\end{exo}

\begin{exo}[fraction massique et pourcentage]
Une solution de masse $\SI{200}{\gram}$ contient $\SI{50}{\gram}$ d'un composé \ce{A} et
$\SI{150}{\gram}$ d'un composé \ce{B}.
\begin{enumerate}[label=\alph*)]
  \item Calcule $w_{\ce{A}}$ et $w_{\ce{B}}$.
  \item Donne le pourcentage massique de chaque composé.
\end{enumerate}
\end{exo}

\begin{exo}[fraction molaire à partir des masses]
Une solution contient $\SI{90}{\gram}$ d'eau \ce{H2O} et $\SI{180}{\gram}$ de glucose \ce{C6H12O6}.
\begin{enumerate}[label=\alph*)]
  \item Calcule la quantité de matière de chaque constituant.
  \item En déduire la fraction molaire de chacun. Vérifie que leur somme vaut $1$.
\end{enumerate}
\end{exo}

\begin{exo}[molalité d'une solution de soude]
On dissout $\SI{4.0}{\gram}$ de \ce{NaOH} dans $\SI{200}{\gram}$ d'eau.
\begin{enumerate}[label=\alph*)]
  \item Quelle quantité de matière de \ce{NaOH} dissout-on~?
  \item Quelle est la molalité de la solution~?
\end{enumerate}
\end{exo}

\end{document}
